\section{Related work}
\label{sec:related}

\subsection{Deep learning for DR grading}
The modern era of automated DR screening started with convolutional networks trained on large collections of graded photographs \citep{gulshan2016}, followed by multi-ethnic validations \citep{ting2017} and regulatory-grade prospective studies \citep{abramoff2018}. Subsequent research on public data has largely consisted of backbone and training-recipe improvements. Residual \citep{he2016resnet}, compound-scaled \citep{tan2019efficientnet}, modernised convolutional \citep{liu2022convnext} and transformer backbones \citep{dosovitskiy2021vit,liu2021swin} have all been applied to grading, typically with high-resolution inputs, heavy augmentation and ordinal or regression objectives. Because retinal images are structurally homogeneous and labelled data are scarce compared with natural images, large self-supervised retinal foundation models such as RETFound \citep{zhou2023retfound} have recently shown that domain-specific pre-training transfers well across diseases and datasets. Our framework is complementary to this trend: it accepts any image embedding (here a frozen ImageNet ConvNeXt-T \citep{liu2022convnext}) and adds lesion-grounded evidence and consistency on top.

A second line of work tackles the specific difficulties of DR grading. The class distribution is heavily skewed towards no-DR and moderate DR, and severe DR is rare. Category-attention modules \citep{he2021cabnet} were proposed to counter this imbalance and were evaluated on DDR. Joint grading of DR and diabetic macular oedema exploits the strong relationship between the two diseases \citep{li2020canet}. Noise-aware and attention-based fusion frameworks \citep{lin2018antinoise} address low-quality inputs. Common to these methods is that they supervise only the grade; the lesions that justify the grade are left implicit.

\subsection{Lesion-aware grading and weak lesion localisation}
Because lesions are the clinical basis of grading, a natural idea is to give the grader access to lesion information. Zoom-in-Net \citep{wang2017zoomin} mines suspicious regions with an attention branch trained only from image-level labels and re-examines them at higher resolution. Image-mining approaches \citep{quellec2017} derive lesion heatmaps from a trained grading network. Two-stage methods first detect lesions and then feed lesion counts or maps to a grader \citep{yang2017twostage}. Lesion-aware transformers \citep{sun2021lat} cast lesion detection and grading in a joint transformer framework, and lesion-based contrastive learning \citep{huang2021lesioncl} uses lesion patches to learn grading-relevant features. These works demonstrate that lesion information improves grading and interpretability. However, to our knowledge, they do not (i) make the grade a function of \emph{zone-resolved} lesion statistics with explicit tokens that can be inspected, (ii) check that the final grade is \emph{consistent} with the lesion evidence according to clinical-rule templates, or (iii) report a rule-violation measure, calibration and conformal coverage jointly. These are the aspects on which \mname\ differs.

\subsection{Retinal lesion segmentation}
Lesion segmentation in fundus images has been driven by encoder--decoder networks descending from U-Net \citep{ronneberger2015}, including nested skip connections \citep{zhou2020unetpp}, attention gates \citep{oktay2018} and atrous multi-scale decoders \citep{chen2018deeplab}. Public pixel-level resources include IDRiD \citep{porwal2020idrid}, e-ophtha \citep{decenciere2013teleophta}, the larger FGADR collection \citep{zhou2021fgadr} and the DDR dataset used here \citep{li2019ddr}. The consistent finding across these benchmarks is that large, high-contrast lesions (hard exudates, large haemorrhages) are segmented reasonably well while microaneurysms remain difficult: they are small, low-contrast, often indistinguishable from vessel fragments, and their annotation is noisy. Overlap measures such as Dice are dominated by large lesions, whereas class-imbalance-robust measures such as the area under the precision--recall curve (AUPR) are more informative for rare, small objects \citep{saito2015pr}. Loss design is therefore important: Tversky-type losses \citep{salehi2017tversky} and their focal variant \citep{abraham2019ftl} weight false negatives more than false positives and are a standard remedy for tiny foreground classes, while focal modulation \citep{lin2017focal} down-weights easy background pixels. We adopt this family and add an explicit full-resolution branch, and we quantify the benefit through an ablation against binary cross-entropy and Dice-based objectives.

\subsection{Lesion detection}
Detection of retinal lesions has mostly reused generic object detectors \citep{ren2017fasterrcnn,tian2019fcos}. These architectures were designed for objects that span tens to hundreds of pixels, with anchors and strides that are poorly matched to lesions that may span only a few pixels, and they require many training boxes, while the DDR detection subset has only 383 training images. Detection from segmentation outputs is an attractive alternative when pixel masks are available: instances are obtained by connected-component analysis and ranked by their mean probability, and a single network serves both tasks. This is the route we take. Following common practice for tiny-object detection, we report average precision at relaxed and standard overlap thresholds \citep{padilla2021detmetrics} together with lesion-level free-response ROC (FROC) sensitivity at fixed numbers of false positives per image, which is the quantity that matters for screening workflows.

\subsection{Multi-task learning and mixed supervision}
Multi-task models for DR combine grading with lesion segmentation to exploit complementary supervision \citep{foo2020multitask}, learn lesion segmentation from image-level labels with weak supervision \citep{playout2019}, or jointly train segmentation and classification in a semi-supervised, collaborative fashion \citep{zhou2019collab}. Attention-based multiple-instance learning \citep{ilse2018attentionmil} provides a general mechanism to localise image-level evidence. In all these designs, sharing is realised through common weights and a summed loss. Our contribution to this line is to make the dependence of the grade on the lesion maps \emph{structural} rather than implicit, and to attach logic-based constraints to that dependence.

\subsection{Ordinal learning, label noise and knowledge-based constraints}
DR grades are ordered, so treating them as nominal categories discards information. Ordinal formulations convert a $K$-class problem to $K\!-\!1$ binary cumulative questions \citep{niu2016ordinal}, with rank-consistent variants that guarantee monotone cumulative probabilities \citep{cao2020coral,shi2023corn}. Inter-grader disagreement in DR concentrates on adjacent grades, which suggests smoothing the target distribution along the grade axis; we adopt a discretised Gaussian target on the cumulative scale. Injecting logical knowledge into neural networks has a long history: logic rules can regularise network outputs through teacher--student distillation \citep{hu2016logic}, semantic-based regularisation with t-norm relaxations \citep{diligenti2017sbr}, or semantic loss functions \citep{xu2018semanticloss}. We use the t-norm relaxation, because it is differentiable, requires no extra labels and gives an interpretable per-rule score. The ICDR/ETDRS grading definitions \citep{wilkinson2003,etdrs1991} supply the templates. To our knowledge such rule templates have not previously been combined with zone-resolved lesion statistics from the same network and an explicit violation-rate metric in DR grading.

\subsection{Calibration and conformal prediction}
Modern networks are over-confident \citep{guo2017calibration}, and for screening the practically important question is what the probability that the true grade lies in the set of grades the system proposes is. Split-conformal prediction \citep{vovk2005,angelopoulos2023} converts any probabilistic classifier into a set-valued predictor with finite-sample marginal coverage under exchangeability, without retraining. Monte-Carlo dropout \citep{gal2016dropout} and deep ensembles estimate epistemic uncertainty but offer no coverage guarantee. We combine temperature scaling for calibration with split-conformal sets, define a clinically motivated referral rule on the resulting sets, and report coverage, set size and referral errors.

\subsection{Positioning}
Table~\ref{tab:lit} summarises how the methods discussed above differ in design. The comparison is qualitative: reported numbers of other works are obtained under different splits, resolutions, backbones and compute budgets, and we have not re-run those systems; therefore we make no numerical claims against them and restrict quantitative comparisons to controlled baselines trained under the same protocol (Section~\ref{sec:setup}).

\begin{table}[!htbp]
\centering
\caption{Design-level comparison with representative prior work (qualitative; $\checkmark$ = addressed explicitly, $\circ$ = partially or implicitly, $-$ = not addressed). Entries reflect our reading of the cited papers' main design and may omit secondary features.}
\label{tab:lit}
\small
\setlength{\tabcolsep}{3.5pt}
\begin{tabular}{@{}p{3.3cm}cccccccc@{}}
\toprule
Method & Grade & Seg. & Det. & Lesion$\to$grade link & Zone-wise evidence & Rule consist. & Ordinal / noise & Conformal \\
\midrule
Zoom-in-Net \citep{wang2017zoomin} & $\checkmark$ & $-$ & $\circ$ & $\circ$ & $-$ & $-$ & $-$ & $-$ \\
Two-stage detect.\ + grade \citep{yang2017twostage} & $\checkmark$ & $-$ & $\checkmark$ & $\checkmark$ & $-$ & $-$ & $-$ & $-$ \\
Foo et al.\ \citep{foo2020multitask} & $\checkmark$ & $\checkmark$ & $-$ & $\circ$ & $-$ & $-$ & $-$ & $-$ \\
Playout et al.\ \citep{playout2019} & $-$ & $\checkmark$ & $-$ & $-$ & $-$ & $-$ & $-$ & $-$ \\
CANet \citep{li2020canet} & $\checkmark$ & $-$ & $-$ & $-$ & $-$ & $-$ & $-$ & $-$ \\
CABNet \citep{he2021cabnet} & $\checkmark$ & $-$ & $-$ & $-$ & $-$ & $-$ & $\circ$ & $-$ \\
Lesion-aware transf.\ \citep{sun2021lat} & $\checkmark$ & $-$ & $\checkmark$ & $\checkmark$ & $-$ & $-$ & $-$ & $-$ \\
Lesion-based contrastive \citep{huang2021lesioncl} & $\checkmark$ & $-$ & $-$ & $\circ$ & $-$ & $-$ & $-$ & $-$ \\
\midrule
\textbf{\mname\ (this work)} & $\checkmark$ & $\checkmark$ & $\checkmark$ & $\checkmark$ & $\checkmark$ & $\checkmark$ & $\checkmark$ & $\checkmark$ \\
\bottomrule
\end{tabular}
\end{table}
